\subsection*{The paradoxes of the theory of sets and the crisis of the foundations}

The first ``paradoxical'' sets appeared in the theory of cardinals and ordinals. In 1897, Burali-Forti noted that there cannot be a set formed of all the ordinals, for such a set would be well-ordered and therefore isomorphic to a proper segment of itself, which is absurd (*). In 1899, Cantor observed (in a letter to Dedekind) that we can no more say that the cardinals form a set, nor speak of ``the set of all sets'' without contradiction (the set of all subsets of the latter ``set'' Ω would be equipotent to a subset of Ω, which is contrary to the inequality m \textless{} 2 \(^{m}\) ) ({[}25{]}, pp.~444-448). Finally, in 1905, by analyzing the proof of this inequality Russell showed that the reasoning which establishes it also proves (without any appeal to the theory of cardinals) that the notion of ``the set of all sets which are not elements of themselves'' is self-contradictory {[}37{]} (†).

It might be thought that such ``antinomies'' would appear only in the peripheral regions of mathematics characterized by the consideration of sets so ``large'' as to be inaccessible to intuition. But other paradoxes were soon to threaten the most classical parts of mathematics. Berry and Russell {[}37{]}, simplifying an argument due to J. Richard {[}34{]}, observed that the set of integers which can be named in less than nineteen syllables is finite, but that it is nevertheless contradictory to define an integer as ``the least integer not nameable in fewer than nineteen syllables'', since this definition contains only eighteen syllables.

Although such arguments, so remote from the current usage of mathematicians, must have appeared to many as a sort of play upon words, yet they indicated the need for a revision of the basis of mathematics in order to eliminate ``paradoxes'' of this nature. But although there was unanimous agreement on the urgency of such a revision, radical divergences of opinion soon arose as to the manner of achieving it. For one group of mathematicians, the ``idealists'' and the ``formalists'' (*), the situation created by the ``paradoxes'' of set theory was very similar to that caused in geometry by the discovery of non-Euclidean geometries or ``pathological'' curves (such as curves with no tangents); it should therefore lead to a similar, but more general conclusion, namely that it is futile to attempt to found any mathematical theory on an appeal (explicit or otherwise) to ``intuition''. This position may be summarized in the words of the principal adversary of the formalist school: ``\ldots{} the view of formalism'', said Brouwer ({[}38 a{]}, p.~83), ``which maintains that human reason does not have at its disposal exact images either of straight lines or of numbers larger than ten, for example\ldots{} It is true that from certain relations among mathematical entities, which we assume as axioms, we deduce other relations according to fixed laws, in the conviction that in this way we derive truths from truths by logical reasoning\ldots{} For the formalist, therefore, mathematical exactness consists merely in the method of developing the series of relations, and is independent of the significance one might want to give to the relations or the entities which they relate''.

For the formalist, therefore, what is required is to provide an axiomatic basis for the theory of sets quite analogous to that for elementary geometry, in which it is not necessary to know what the ``things'' are that are called ``sets'', nor what the relation \(x \in y\) means, but in which the conditions imposed on this relation are enumerated; naturally, this should be done in such a way as to include, so far as possible, all the results of Cantor's theory, while rendering impossible the existence of ``paradoxical'' sets. The first example of such an axiomatization was given by Zermelo in 1908 {[}36 c{]}; he avoided sets which are ``too big'' by introducing a ``selection axiom'' (Aussonderungsaxiom), according to which a property \(P\{x\}\) determines a set formed by the elements which possess this property only if \(P\{x\}\) already implies a relation of the form \(x \in A (\dagger)\) . But the elimination of paradoxes analogous to ``Richard's paradox'' could be achieved only by restricting the sense attached to the notion of ``property''.

In this connection, Zermelo was content to describe in extremely vague terms a type of properties which he called ``definite'' and to indicate that the application of the selection axiom must be limited to properties of this type. This point was made precise by Skolem {[}39{]} and Fraenkel {[}41{]} (*); as they observed, its elucidation demands a completely formalized system (such as that described in Chapters I and II of this book), in which the notions of ``property'' and ``relation'' have lost all ``meaning'' and have become simply designations for assemblies formed in accordance with explicit laws. This, of course, necessitates that the rules of logic used be incorporated in the system, which was not yet the case in the Zermelo-Fraenkel system; apart from this, it is essentially their system which we have described in Chapters I and II.

Other axiomatizations of the theory of sets were subsequently proposed. We shall consider principally that due to von Neumann {[}43 a and b{]}, which comes closer than the Zermelo-Fraenkel system to Cantor's primitive conception; the latter, in order to avoid ``paradoxical'' sets, had already proposed (in his correspondence with Dedekind; {[}25{]}, pp.~445-448) to distinguish two kinds of sets, ``multiplicities'' (Vielheiten) and ``sets'' (Mengen) in the strict sense, the latter being characterized by the fact that they could be thought of as single objects. It is this idea which von Neumann made precise by distinguishing between two types of objects, ``sets'' and ``classes''. In his system (almost completely formalized) the classes are distinguished from sets by the fact that they cannot be placed on the left of the sign \(\in\) . One of the advantages of such a system is that it rehabilitates the notion of ``universal class'' (which, of course, is not a set) used by the logicians of the nineteenth century. We should add that von Neumann's system avoids (for set theory) the introduction of axiom schemes, which are replaced by suitable axioms, thereby making the logical study of the system easier. Variants of von Neumann's system have been given by Bernays and Gödel {[}44 b{]}.

These systems seem to have succeeded in eliminating the paradoxes, but at the cost of restrictions which cannot but appear highly arbitrary. In favour of the Zermelo-Fraenkel system it can be said that it limits itself to formulating prohibitions which do no more than sanction current practice in the applications of the notion of set to various mathematical theories. The systems of von Neumann and Gödel are more remote from usual conceptions. On the other hand, we cannot exclude the possibility that it may be easier to insert the basis of some mathematical theories into the framework of such systems than into the more rigid framework of the Zermelo-Fraenkel system.

In any case, it cannot be said that any of these solutions gives the impression of finality. They satisfy the formalists, because the formalists refuse to take into consideration the individual psychological reactions of every mathematician; they consider that a formalized language has done its duty if it is capable of transcribing mathematical arguments in a form which contains no ambiguities, and hence of serving as a vehicle for mathematical thought. It is everyone's right, they would say, to think what they like about the ``nature'' of mathematical entities or the ``truth'' of the theorems they use, provided that their reasoning can be transcribed into the common language (*).

In other words, from a philosophical point of view, the attitude of the formalists is characterized by an indifference to the problem posed by the ``paradoxes'', resulting from the abandonment of the Platonic position which claimed to attribute to mathematical notions an intellectual ``content'' common to all mathematicians. Many mathematicians rebelled against this break with tradition. Russell, for example, sought to avoid the paradoxes by analyzing their structure more deeply. Taking up an idea first advanced by J. Richard (in the article {[}34{]} in which he presented his ``paradox'') and later developed by H. Poincaré {[}33 c{]}, Russell and Whitehead observed that the definitions of the paradoxical sets all violate the following principle, called the ``vicious circle principle'': ``whatever involves all of a collection must not be one of the collection'' ({[}37{]}, vol.~I, p.~40). This statement served as a basis for the Principia, and the ``theory of types'' was developed in order to accommodate it. Like that of Frege which inspired it, the logic of Russell and Whitehead (much vaster than the mathematical logic used in Chapter I of this Book) contains ``propositional variables''. The theory of types proceeds to a classification of these variables, in broad outline as follows.

Starting from an undefined ``domain of individuals'', which may be qualified as ``objects of order 0'', relations in which the variables (free or bouud) are individuals are called ``first-order objects''; and in general, relations in which the variables are objects of order \(\leqslant n\) (and at least one of them is of order n) are called ``objects of order \(n + 1\) '' (*). A set of objects of order n can therefore be defined only by a relation of order \(n + 1\) , and this condition allows the elimination of paradoxical sets without any difficulty ( \(\dagger\) ). But the principle of the ``hierarchy of types'' is so restrictive that, if strictly adhered to, it leads to a mathematics of inextricable complexity (**). To escape this consequence, Russell and Whitehead were obliged to introduce an ``axiom of reducibility'' which asserts the existence, for every relation between ``individuals'', of an equivalent relation of the first order. This condition is as arbitrary as the axioms of the formalists, and reduces considerably the interest of the construction of the Principia. The system of Russell and Whitehead has been more popular with logicians than with mathematicians. Moreover, it is not completely formalized ( \(\dagger\dagger\) ), so that there are numerous obscurities of detail. Various efforts have been made to simplify and clarify it (Ramsey, Chwistek, Quine, Rosser); these authors tend to use more and more completely formalized languages, and replace the rules of the Principia (which still contain a certain intuitive substratum) by restrictions which take account only of the writing down of the assemblies considered; not only do these rules thus appear as gratuitous as the prohibitions formulated in the systems of Zermelo-Fraenkel or von Neumann, but also, being more remote from mathematical usage, they have in many cases led to unacceptable consequences which the author had not foreseen (such as Burali-Forti's paradox, or the negation of the axiom of choice).

For the mathematicians of these schools, the essential aim was to avoid renouncing any part of the heritage of the past. ``No one shall expel us'', said Hilbert ({[}30{]}, p.~274), ``from the paradise Cantor has created for us''. To achieve this aim, they were prepared to accept limitations on mathematical reasoning which, though not of vital importance because they conformed with mathematical usage, did not appear to be imposed by our mental habits and the intuitive notion of a set. For them, anything was preferable to the intrusion of psychology in the criteria of validity in mathematics; and rather than ``take into account the properties of our brains'', as Hadamard said ({[}35{]}, p.~270), they resigned themselves to the imposition of largely arbitrary boundaries, provided that all of classical mathematics was contained in them and that they did not threaten to impede later progress.

Quite different was the attitude of the mathematicians who adhered to the following doctrine. Whereas the formalists agreed to abandon control by the ``eyes of the mind'' so far as mathematical reasoning was concerned, the mathematicians who have been labeled as ``empiricists'', ``realists'', or ``intuitionists'' refused to consent to this abdication; they insisted on a sort of inner certainty guaranteeing the ``existence'' of the mathematical objects they studied. They had no serious objection to the renunciation of spatial intuition, because arithmetical ``models'' allowed them to shelter behind the intuitive notion of integers. But irreconcilable opposition arose over the question of reducing the notion of integer to that of set (which is intuitively far less clear) and then of erecting barriers, devoid of intuitive foundations, to the manipulation of sets. The first person to voice this opposition (and the most influential, by reason of the authority of his genius) was H. Poincaré. Having accepted not only the axiomatic point of view toward geometry and the arithmetization of analysis but also a good part of Cantor's theory, he refused to consider that arithmetic, too, might be amenable to axiomatic treatment. The principle of mathematical induction, in particular, was in his eyes a fundamental intuition of our intellects and could not be regarded purely as a convention (*) {[}33 c{]}. He was in principle against formalized languages, whose usefulness he contested, and was prone to confuse the notion of integers in formalized mathematics with the use of integers in the theory of proof (which was then emerging and of which we shall speak later). No doubt it was not easy at that time to make this distinction as precisely as we do today --- after more than fifty years of study and discussion --- although the distinction had been grasped clearly by men such as Hilbert and Russell.

Criticisms of this nature multiplied after the introduction of the axiom of choice by Zermelo in 1904 {[}36 a{]}. Its use in many previous proofs in analysis and set theory had hitherto gone almost unnoticed (*), and it was by following up an idea suggested by Erhard Schmidt that Zermelo, explicitly stating this axiom (and in the two forms given in the Summary of Results, §4, no. 10), deduced by an ingenious method (which we reproduced in essence in Chapter III, §2) a satisfactory proof of the existence of a well-ordering on any set. It appears that, coming as it did at the same time as the ``paradoxes'', this new method of reasoning, with its unfamiliar appearance, caused confusion among many mathematicians; one need only read the strange misconceptions which arose in this connection, in the following volume of the Mathematische Annalen, by authors as familiar with Cantor's methods as Schoenflies and F. Bernstein. The criticisms of E. Borel, published in the same volume, have more substance and clearly reflect the point of view of Poincaré on integers; they were developed and discussed in an exchange of letters between E. Borel, Baire, Hadamard, and Lebesgue, which has become a classic of the French mathematical tradition {[}35{]}. Borel began by denying the validity of the axiom of choice because in general it involves an uncountable infinity of choices, which cannot be intuitively imagined. But Hadamard and Lebesgue observed that a countable infinity of successive arbitrary choices is no more intuitive, because it involves an infinite number of operations, which it is impossible to conceive of being actually carried out. In the view of Lebesgue, who enlarged the debate, everything depended on what was meant by the assertion that a mathematical object ``exists''; for him, it is necessary to ``name'' explicitly a property (a ``functional'' property, we should say) which defines the object uniquely. A function such as that used by Zermelo in his proof is what Lebesgue called a ``law'' of choice. If, he continued, this requirement is not satisfied, and we merely ``consider'' this function instead of ``naming'' it, can we be sure that throughout our reasoning we are always considering the same function? ({[}35{]}, p.~267) This consideration led Lebesgue to new doubts: even the choice of a single element in a set appears to raise difficulties; it is necessary to be certain that such an element ``exists'', i.e., that at least one element of the set can be ``named'' (†). Can we therefore speak of the ``existence'' of a set if we cannot ``name'' every element of it? Baire did not hesitate to deny the ``existence'' of the set of subsets of a given infinite set (loc. cit., pp.~263-264); in vain Hadamard observed that these restrictions would lead to renouncing even the right to speak of the set of real numbers, and E. Borel in the end came to this conclusion. Apart from the fact that countable sets seem to have been granted the franchise, we are almost back to the classical position of the opponents of the ``actual infinite''.

None of these objections was very systematic; it was reserved to Brouwer and his school to undertake a complete recasting of mathematics on similar but even more radical lines. We shall not here presume to summarize a doctrine so complex as intuitionism, which is as much of philosophy as mathematics, and we shall do no more than indicate some of its most striking features, referring for more details to the works of Brouwer himself {[}38{]} and Heyting's exposition {[}46{]}. In Brouwer's view, mathematics is identical with the ``exact'' part of our thought, founded on the primary intuition of the sequence of natural numbers, and cannot be translated into a formal system without mutilation. Moreover, it is ``exact'' only in the minds of mathematicians, and it is illusory to hope to construct an instrument of communication between mathematicians which is free of all the imperfections and ambiguities of languages; the most that can be hoped for is to arouse in the interlocutor a favourable state of mind for more or less vague descriptions ({[}46{]}, pp.~11-13). Intuitionist mathematics attaches hardly more importance to logic than to language; a proof is conclusive not by virtue of the rules of logic, fixed once and for all, but by reason of the ``immediate self-evidence'' of each of its steps. This ``self-evidence'' is moreover to be interpreted in an even more restrictive sense than that of E. Borel and his supporters; thus in intuitionist mathematics we may not assert that a relation of the form ``R or (not R)'' is true (the principle of the excluded middle) unless, for every system of values given to the variables appearing in R, we can prove that one of the two propositions R, ``not R'' is true. For example, from the equation \(ab = 0\) between two real numbers, we may not infer '' \(a = 0\) or \(b = 0\) ``, because it is easy to construct explicit examples of real numbers \(a\) , \(b\) such that we have \(ab = 0\) without at present being able to prove either one of the two propositions \(a = 0\) , \(b = 0\) ({[}36{]}, p.~21).

It is hardly surprising that the intuitionist mathematicians, starting from such principles, were led to results quite different from the classical theorems. A large number of the latter disappeared, for example most of the ``existence'' theorems of analysis (such as the theorems of Bolzano and Weierstrass for real functions); if a function of a real variable ``exists'' in the intuitionist sense, it is ipso facto continuous, and a monotone bounded sequence of real numbers does not necessarily have a limit. Furthermore, many of the classical notions ramify into several fundamentally distinct notions in intuitionist mathematics. Thus there are two notions of convergence (for a sequence of real numbers) and eight notions of countability. It goes without saying that transfinite induction and its applications to modern analysis are, like most of Cantor's theory, condemned without appeal.

According to Brouwer, it is only in this way that mathematical propositions can acquire a ``content''; formalist arguments which go beyond what is admitted by intuitionism are judged to be worthless, because it is impossible to give them a ``meaning'' to which the intuitive notion of ``truth'' could be applied. Clearly, judgments of this sort can only rest on a previous notion of ``truth'', which is of a psychological or metaphysical nature; that is to say, in practice, that they are beyond discussion.

Certainly, the vigorous attacks from the intuitionist camp have from time to time obliged not only the avant-garde mathematical schools, but even the partisans of traditional mathematics, to be on the defensive. A well-known mathematician has admitted to being impressed by these attacks to the point that he voluntarily restricted his work to those branches of mathematics considered to be ``certain''. But such cases must have been rather uncommon. The intuitionist school, whose memory will undoubtedly survive only as a historical curiosity, has at least rendered the service of having obliged its opponents, that is to say the vast majority of mathematicians, to clarify their own positions and to become more consciously aware of the reasons (whether logical or sentimental) for their confidence in mathematics.

